\documentclass[10pt,letterpaper]{article}

\usepackage[letterpaper,margin=0.58in]{geometry}
\usepackage[T1]{fontenc}
\usepackage[utf8]{inputenc}
\usepackage{lmodern}
\usepackage{microtype}
\usepackage{enumitem}
\usepackage{tabularx}
\usepackage{titlesec}
\usepackage{xcolor}
\usepackage[hidelinks]{hyperref}

\definecolor{burgundy}{HTML}{6E2639}
\definecolor{indigo}{HTML}{31365A}
\definecolor{muted}{HTML}{555555}

\pagestyle{empty}
\setlength{\parindent}{0pt}
\setlength{\parskip}{0pt}
\setlist[itemize]{leftmargin=1.05em,itemsep=1.5pt,topsep=2pt,parsep=0pt,partopsep=0pt}
\titleformat{\section}{\large\bfseries\color{burgundy}}{}{0pt}{}[\vspace{-3pt}\color{burgundy}\titlerule]
\titlespacing*{\section}{0pt}{7pt}{4pt}
\hypersetup{
  colorlinks=true,
  urlcolor=indigo,
  pdftitle={Samantha Lune Magid - Curriculum Vitae},
  pdfauthor={Samantha Lune Magid},
  pdfsubject={Academic curriculum vitae}
}

\newcommand{\entry}[4]{%
  \begin{tabularx}{\textwidth}{@{}Xr@{}}
    \textbf{#1} & \textbf{#2} \\
    \textit{#3} & \textit{#4}
  \end{tabularx}\vspace{-2pt}
}

\newcommand{\skill}[2]{\textbf{#1:} #2\par\vspace{2pt}}

\begin{document}

\begin{center}
  {\LARGE\bfseries\color{burgundy} Samantha Lune Magid}\par
  \vspace{2pt}
  Applied and Computational Mathematician\par
  \vspace{3pt}
  \href{mailto:samantha.magid@uvm.edu}{samantha.magid@uvm.edu}
  \enspace$\vert$\enspace Burlington, Vermont
  \enspace$\vert$\enspace \href{https://magid.systems}{magid.systems}
  \enspace$\vert$\enspace \href{https://github.com/crabwife}{github.com/crabwife}
  \enspace$\vert$\enspace \href{https://orcid.org/0009-0004-0358-4870}{ORCID}
\end{center}

\section*{Research Interests}
Complex and adaptive systems; stochastic processes and history-dependent dynamics; science of science and metaresearch; scientific productivity and career trajectories; discrete-time count and sequential processes; emergent behavior in social and ecological systems; geometric and latent representations; causal inference and quasi-experimental identification; interpretable generative modeling; estuarine and intertidal systems; coupled air-sea climate systems; ethical limits of mathematical modeling.

\section*{Education}
\entry{University of Vermont}{Expected Spring 2028}{M.S. Complex Systems and Data Science, Vermont Complex Systems Institute}{Burlington, VT}
\begin{itemize}
  \item UVM Graduate Merit Scholar.
  \item Research areas: stochastic scientific productivity, science of science, metaresearch, and computational social science.
\end{itemize}

\entry{Boston University}{B.S.}{B.S. Data Science, Applied and Mathematical Focus; Minor in Environmental Science}{Boston, MA}
\begin{itemize}
  \item Lara Vincent Undergraduate Research Scholar and Award recipient, October 2022.
\end{itemize}

\section*{Research Experience}
\entry{Graduate Research Assistant}{May 2026--Present}{Science and Humanity Lab, Vermont Complex Systems Institute}{Burlington, VT}
\begin{itemize}
  \item Develop scholar-agnostic stochastic models of scientific productivity to test whether persistent career differences can emerge from dynamic processes rather than fixed individual traits.
  \item Fit and simulate autoregressive, hurdle, self-exciting, and history-dependent models using longitudinal publication data from DBLP and the Academic Analytics Research Center.
  \item Identified distinct extensive and intensive margins of productivity, geometric memory decay, declining restart hazards after inactivity, and persistent but dynamically generated rank structure.
  \item Design distributional, predictive, and trajectory-level evaluations using held-out simulation, bootstrap confidence intervals, cross-validation, rank-mixing analyses, goodness-of-fit tests, and adversarial null models.
  \item Built reproducible Python pipelines for data preparation, model fitting, simulation, diagnostics, and publication-quality visualization across multiple linked research repositories.
  \item Contribute to manuscript development, figures, technical reports, conference materials, and a multi-field extension of prior work on scientific productivity as a stochastic process.
\end{itemize}

\section*{Manuscripts and Presentations}
\textbf{Zhang, S.; \textbf{Magid, S.}; Kumar, A.; Larremore, D. B.; Clauset, A.}
``Stochastic dynamics can drive large differences in researcher productivity.'' \textit{Manuscript in preparation.}

\vspace{2pt}
Related findings were presented by Sam Zhang as a lightning talk at the \textit{5th International Conference on the Science of Science and Innovation (ICSSI 2026)}, Boulder, Colorado, June 29, 2026.

\vspace{3pt}
\textbf{The Productivity Garden.} Independent paper in progress developing a stochastic project-pipeline model of publication timing, attrition, memory, and clustered output.

\section*{Research Software}
\textbf{Selected repository:} \href{https://github.com/crabwife/spaarw_followup}{spaarw\_followup} --- model lineage, memory tests, trajectory simulation, empirical diagnostics, and reproducible figures for the scientific-productivity project. Additional public code includes \href{https://github.com/crabwife/seven_models_shockley}{seven\_models\_shockley} and \href{https://github.com/crabwife/AARC-DBLP-comparison}{AARC-DBLP-comparison}.

\newpage

\section*{Laboratory and Teaching Experience}
\entry{Graduate Laboratory Supervisor / Advanced Lab Professional II}{July 2026--Present}{Human Cell Biology Program, University of Vermont}{Burlington, VT}
\begin{itemize}
  \item Run the hands-on operations of the Human Cell Biology preparation laboratory, coordinating materials, protocols, equipment, and readiness for instructional labs.
  \item Supervise preparation-lab technicians and coordinate with faculty and staff to support safe, consistent, and technically reliable laboratory instruction.
\end{itemize}

\section*{Undergraduate Research}
\entry{Undergraduate Principal Investigator}{September 2022--March 2023}{\textit{Uca spp.} Bioturbation Survey, Boston University Marine Program}{Boston, MA}
\begin{itemize}
  \item Proposed and led an independent study of fiddler-crab overpopulation as a complex biophysical system, examining how stochastic bioturbation feedback cycles contribute to salt-marsh erosion.
  \item Designed and fabricated a custom tidal-forcing generator for controlled simulation of periodic intertidal conditions without disrupting field sites.
  \item Coordinated with marine scientists and field experts, completed the project's first phase, and produced reusable tools for future estuarine-modeling research.
\end{itemize}

\section*{Mathematical and Computational Skills}
\skill{Stochastic and complex systems}{Stochastic and Markov processes; discrete-time count and sequential models; autoregressive, hurdle, survival, hazard, and history-dependent models; geometric memory kernels; feedback-aware dynamical systems; emergence and nonlinearity; stochastic biological systems; latent state-space models.}
\skill{Statistical and computational methods}{Monte Carlo simulation; bootstrap inference; held-out simulation; cross-validation; goodness-of-fit testing; time-series analysis; low-dimensional approximations; rank-persistence and transition analysis; robustness and sensitivity analysis.}
\skill{Optimization, learning, and causal methods}{Discrete and continuous optimization; kernel methods; geometric representation; neural-network architectures; black-box optimization; simulated annealing; causal inference; quasi-experimental identification; adversarial and game-theoretic models.}
\skill{Programming}{Python, R, SQL, Rust.}
\skill{Python scientific stack}{pandas, NumPy, SciPy, statsmodels, scikit-learn, JAX, NumPyro, PyTorch, TensorFlow, XGBoost, Matplotlib, seaborn.}
\skill{Research and visualization tools}{Git, GitHub, Jupyter, Quarto, LaTeX, QGIS, ArcGIS, NetCDF, MySQL, Power BI, Tableau.}
\skill{Laboratory and field}{Preparation-lab supervision; wet-lab techniques; animal handling; laboratory safety; IACUC certification; experimental-system design.}

\section*{Affiliations}
Member, American Statistical Association (AmStat) \enspace$\vert$\enspace Member, Society for Industrial and Applied Mathematics (SIAM) \enspace$\vert$\enspace Vermont Complex Systems Institute \enspace$\vert$\enspace Science and Humanity Lab

\vfill
\begin{center}
  {\small\color{muted} Updated September 2026}
\end{center}

\end{document}
